%% appendices/appendixA.tex
%% ============================================================================
%% APPENDIX A: SOURCE CODE
%% ============================================================================

\chapter{Source Code}
\label{app:code}

This appendix presents the core Python classes of the HippoCortex Stage 1 continual learning implementation.

\section{StatsBuffer Module}
\label{sec:app_stats_buffer}

The \texttt{StatsBuffer} class implements the per-task compact memory storage. It computes and stores the mean and variance of hidden states for completed tasks, ensuring that memory consumption scales strictly as $\mathcal{O}(T \times d_{\text{model}})$.

\begin{lstlisting}[language=Python, caption={StatsBuffer Implementation}, label={lst:app_stats_buffer_code}]
import torch
from torch import Tensor

class StatsBuffer:
    def __init__(self) -> None:
        self._store: dict[int, tuple[Tensor, Tensor]] = {}

    def update(self, task_id: int, hidden_states: Tensor) -> None:
        """
        Compute and store (mean, variance) for task_id.
        Args:
            task_id: Integer task index (0-based).
            hidden_states: Shape (N, d_model) - hidden states collected
                           across the training set of the completed task.
        """
        mu = hidden_states.mean(dim=0)
        sigma2 = hidden_states.var(dim=0, unbiased=False)
        self._store[task_id] = (mu, sigma2)

    def get_stats(self, task_id: int) -> tuple[Tensor, Tensor]:
        """
        Return stored (mean, variance) for task_id.
        Returns:
            mu:     shape (d_model,)
            sigma2: shape (d_model,)
        """
        if task_id not in self._store:
            raise KeyError(f"Task {task_id} has not been stored yet.")
        return self._store[task_id]

    def task_ids(self) -> list[int]:
        """Return sorted list of all stored task ids."""
        return sorted(self._store.keys())

    def memory_bytes(self) -> int:
        """
        Return exact byte count of all stored tensors.
        Satisfies: bytes == len(task_ids) * d_model * 2 * element_size.
        """
        if not self._store:
            return 0
        first_mu = next(iter(self._store.values()))[0]
        d_model = first_mu.shape[0]
        element_size = first_mu.element_size()
        return len(self._store) * d_model * 2 * element_size
\end{lstlisting}

\section{NullSpaceProjector Module}
\label{sec:app_null_space_projector}

The \texttt{NullSpaceProjector} class implements the gradient orthogonalisation logic. It performs incremental Singular Value Decomposition (SVD) and QR orthonormalisation to maintain a projection basis, projecting gradients into the null-space of prior task activations.

\begin{lstlisting}[language=Python, caption={NullSpaceProjector Implementation}, label={lst:app_null_space_code}]
import torch
from torch import Tensor

class NullSpaceProjector:
    def __init__(self, rank_budget: int = 200) -> None:
        """
        Args:
            rank_budget: Maximum number of basis vectors to retain.
        """
        self.rank_budget = rank_budget
        self._U: Tensor | None = None  # shape (d_model, k), k <= rank_budget

    def update(self, hidden_states: Tensor) -> None:
        """
        Extend the basis U with new task's feature directions.
        Args:
            hidden_states: Shape (N, d_model) - completed task states.
        """
        _, S, Vt = torch.linalg.svd(hidden_states, full_matrices=False)
        var = S ** 2
        cumulative_ratio = torch.cumsum(var, dim=0) / var.sum()
        r = int((cumulative_ratio < 0.99).sum().item()) + 1
        r = min(r, Vt.shape[0])
        new_dirs = Vt[:r].T
        
        if self._U is None:
            combined = new_dirs
        else:
            combined = torch.cat([self._U, new_dirs], dim=1)
            
        Q, _ = torch.linalg.qr(combined)
        k_new = min(Q.shape[1], self.rank_budget)
        self._U = Q[:, :k_new]

    def project(self, grad: Tensor) -> Tensor:
        """
        Apply the null-space projection P = I - U @ U.T to a gradient tensor.
        Args:
            grad: Gradient of shape (..., d_model).
        """
        if self._U is None:
            return grad
        return grad - (grad @ self._U) @ self._U.T

    @property
    def current_rank(self) -> int:
        """Number of basis vectors currently stored."""
        return 0 if self._U is None else self._U.shape[1]
\end{lstlisting}
